\chapter{2026 年 818 工程流体力学试题（考生回忆版）}

\begin{tishi}
\textbf{使用说明：}本套为考生回忆版，题干文字与原卷可能有出入，个别题目的数值与顺序不保证准确，
但\textbf{考点分布、题型配比、命题方式}具有很高的参考价值。回忆版中明确点出的
"书上例题 4.12""书上例题（管嘴）"等，请对照《教材精讲》第 4、5 章相应考点复习。
\end{tishi}

\section{一、单项选择题（10 小题，每题 3 分，共 30 分）}

\begin{tishi}
回忆者注：「\textbf{和往年真题还有题库类似}」。
——即\textbf{直接来自历年真题与题库的原题或改编题}，这是本套题最重要的一条信息：
选择题不需要去猜新考点，把 2014--2022 年真题与题库的选择题做熟即可覆盖。
\end{tishi}

\section{二、填空题（10 空，每空 3 分，共 30 分）}

\begin{tishi}
回忆者注：「\textbf{同理}」（同样与往年真题、题库类似）。填空题以名词定义与公式含义为主，
重点复习《教材精讲》中每个考点的"是什么"与"公式含义"两行。
\end{tishi}

\section{三、简答题（3 小题，每题 4 分，共 12 分）}

\begin{ti}
  \item \textbf{水击现象}：什么是水击？产生的原因是什么？有哪些危害与防护措施？
  \item \textbf{理想流体}：什么是理想流体？引入理想流体概念的意义是什么？
  \item \textbf{层流与雷诺数的关系}：层流与湍流如何判别？雷诺数的物理意义是什么？
\end{ti}

\section{四、计算题（5 题，共 58 分）}

\begin{ti}
  \item \textbf{牛顿流体（斜板静压强）}：一块平板（左高右低）斜置于流体中，求板上的静压强 / 切应力分布。
  \item \textbf{文丘里管流量的证明}：证明文丘里流量计的流量公式 $Q=\mu A\sqrt{2g\Delta h/(\ldots)}$。
  \item \textbf{动量方程（教材例题 4.12 原题）}：教材上动量方程的例题，务必把 4.12 题的原解法背熟。
  \item \textbf{证明势函数是否存在}：给定速度场，判断/证明该流动是否存在速度势函数 $\varphi$。
  \item \textbf{管嘴体积流量（教材例题）}：教材上管嘴出流的例题，求体积流量。
\end{ti}

\vspace{4mm}
\begin{jielun}
\textbf{从这一套回忆版能读出的命题趋势（用于指导复习优先顺序）：}
\begin{enumerate}
  \item 选择、填空\textbf{几乎全部来自"真题＋题库"的重复}——重复率高的判断再次被验证；
  \item 简答考的是\textbf{第 5 章（水击）、第 1 章（理想流体）、第 5 章（层流与雷诺数）}的定义型内容；
  \item 计算题 5 题中，\textbf{2 题直接是教材例题}（动量方程 4.12、管嘴），
        说明\textbf{教材例题与课后习题的改编是命题的主要来源}；
  \item 计算题覆盖第 1 章（牛顿流体）、第 4 章（文丘里、动量方程）、第 7 章（势函数）、第 5 章（管嘴），
        \textbf{与考情分析"第四、五章是计算题重点"的结论一致}。
\end{enumerate}
\end{jielun}
